% c12 原书 p.149–160（PDF p.156–167）
\subsection{谱序列}
\label{sec:2.2}

设 $\cC$ 是阿贝尔范畴。$\cC$ 中的\emph{谱序列}\idx{puxulie}{谱序列} $E$ 由以下数据组成：
\begin{enumerate}[label=(\alph*),leftmargin=2.8em]
  \item 一族 $\cC$ 中的对象 $E_r^{pq}$，其中 $p,q,r\in\bZ$ 且 $r\ge 2$。
  \item 一族态射 $d_r^{pq}\colon E_r^{pq}\to E_r^{p+r,q-r+1}$，满足 $d_r^{p+r,q-r+1}d_r^{pq}=0$。令
  \[
  Z_{r+1}(E_r^{pq})=\Ker d_r^{pq},\qquad B_{r+1}(E_r^{pq})=\Image d_r^{p-r,q+r-1}.
  \]
  则有 $B_{r+1}(E_r^{pq})\subset Z_{r+1}(E_r^{pq})\subset E_r^{pq}$。
  \item 一族同构
  \[
  \alpha_r^{pq}\colon Z_{r+1}(E_r^{pq})/B_{r+1}(E_r^{pq})\cong E_{r+1}^{pq}.
  \]
\end{enumerate}

对每个 $k\ge r+1$，我们用对 $k-(r+1)$ 的归纳法定义 $E_r^{pq}$ 的一些子对象 $B_k(E_r^{pq})$ 与 $Z_k(E_r^{pq})$。当 $k-(r+1)=0$ 时，$B_k$ 与 $Z_k$ 已由 (b) 定义。设对一切满足 $0\le i-(s+1)<k-(r+1)$ 的 $i,s$，$B_i(E_s^{pq})$ 与 $Z_i(E_s^{pq})$ 都已定义；特别地，$E_{r+1}^{pq}$ 的子对象 $B_k(E_{r+1}^{pq})$ 与 $Z_k(E_{r+1}^{pq})$ 已定义。通过同构
\[
\alpha_r^{pq}\colon Z_{r+1}(E_r^{pq})/B_{r+1}(E_r^{pq})\cong E_{r+1}^{pq},
\]
这两个子对象等同于 $Z_{r+1}(E_r^{pq})/B_{r+1}(E_r^{pq})$ 的某两个子对象。定义 $B_k(E_r^{pq})$ 与 $Z_k(E_r^{pq})$ 为它们在典范满射
\[
Z_{r+1}(E_r^{pq})\to Z_{r+1}(E_r^{pq})/B_{r+1}(E_r^{pq})
\]
下的原像。此时有 $B_k\subset Z_k$，且
\[
Z_k(E_r^{pq})/B_k(E_r^{pq})\cong Z_k(E_{r+1}^{pq})/B_k(E_{r+1}^{pq})\cong\cdots\cong Z_k(E_{k-1}^{pq})/B_k(E_{k-1}^{pq})\cong E_k^{pq}.
\]

若令 $B_r(E_r^{pq})=0$、$Z_r(E_r^{pq})=E_r^{pq}$，则
\[
0=B_r(E_r^{pq})\subset B_{r+1}(E_r^{pq})\subset B_{r+2}(E_r^{pq})\subset\cdots,
\]
且
\[
\cdots\subset Z_{r+2}(E_r^{pq})\subset Z_{r+1}(E_r^{pq})\subset Z_r(E_r^{pq})=E_r^{pq},
\]
并且第一行中的每个对象都是第二行中相应对象的子对象。

谱序列的定义还包含以下数据：
\begin{enumerate}[label=(\alph*),leftmargin=2.8em,start=4]
  \item $E_2^{pq}$ 的两个子对象 $B_\infty(E_2^{pq})$ 与 $Z_\infty(E_2^{pq})$，使对一切 $k\ge 2$ 有
  \[
  B_k(E_2^{pq})\subset B_\infty(E_2^{pq})\subset Z_\infty(E_2^{pq})\subset Z_k(E_2^{pq}).
  \]
  令 $E_\infty^{pq}=Z_\infty(E_2^{pq})/B_\infty(E_2^{pq})$。
  \item 一族 $\cC$ 中的对象 $H^n$（$n\in\bZ$），以及每个 $H^n$ 上的一个\emph{递降滤过}\idx{dijianglvguo}{递降滤过}：
  \[
  H^n\supset\cdots\supset F^pH^n\supset F^{p+1}H^n\supset\cdots.
  \]
  \item 一族同构
  \[
  \beta^{pq}\colon E_\infty^{pq}\cong F^pH^{p+q}/F^{p+1}H^{p+q}.
  \]
\end{enumerate}

至此，谱序列 $E$ 的定义完成。我们常把它记为 $E_2^{pq}\Rightarrow H^{p+q}$。

设 $E$ 与 $E'$ 是两个谱序列。谱序列的态射 $u\colon E\to E'$ 由一族态射 $u_r^{pq}\colon E_r^{pq}\to E'^{pq}_r$ 与一族态射 $u^n\colon H^n\to H'^n$ 组成，要求它们满足以下条件：
\begin{enumerate}[label=(\roman*),leftmargin=2.8em]
  \item $u^n(F^pH^n)\subset F^pH'^n$，且下图交换：
  \[\begin{tikzcd}[column sep=3.2em,row sep=2.2em]
  E_r^{pq} \arrow[r,"d_r^{pq}"] \arrow[d,"u_r^{pq}"'] & E_r^{p+r,q-r+1} \arrow[d,"u_r^{p+r,q-r+1}"] \\
  E'^{pq}_r \arrow[r,"d'^{pq}_r"'] & E'^{p+r,q-r+1}_r
  \end{tikzcd}\]
  于是 $u^n$ 诱导态射
  \[
  \ol u^n\colon F^pH^n/F^{p+1}H^n\to F^pH'^n/F^{p+1}H'^n,
  \]
  $u_r^{pq}$ 诱导态射
  \[
  \ol u_r^{pq}\colon Z_{r+1}(E_r^{pq})/B_{r+1}(E_r^{pq})\to Z_{r+1}(E'^{pq}_r)/B_{r+1}(E'^{pq}_r).
  \]
  \item 下图交换：
  \[\begin{tikzcd}[column sep=4.6em,row sep=2.4em]
  Z_{r+1}(E_r^{pq})/B_{r+1}(E_r^{pq}) \arrow[r,"\alpha_r^{pq}"] \arrow[d,"\ol u_r^{pq}"'] & E_{r+1}^{pq} \arrow[d,"u_{r+1}^{pq}"] \\
  Z_{r+1}(E'^{pq}_r)/B_{r+1}(E'^{pq}_r) \arrow[r,"\alpha'^{pq}_r"'] & E'^{pq}_{r+1}
  \end{tikzcd}\]
  于是对一切 $k\ge r\ge 2$ 有
  \[
  u_r^{pq}(B_k(E_r^{pq}))\subset B_k(E'^{pq}_r),\qquad u_r^{pq}(Z_k(E_r^{pq}))\subset Z_k(E'^{pq}_r).
  \]
  \item $u_2^{pq}(B_\infty(E_2^{pq}))\subset B_\infty(E'^{pq}_2)$，$u_2^{pq}(Z_\infty(E_2^{pq}))\subset Z_\infty(E'^{pq}_2)$。于是 $u_2^{pq}$ 诱导态射
  \[
  u_\infty^{pq}\colon E_\infty^{pq}\to E'^{pq}_\infty.
  \]
  \item 下图交换：
  \[\begin{tikzcd}[column sep=4.6em,row sep=2.4em]
  E_\infty^{pq} \arrow[r,"\beta^{pq}"] \arrow[d,"u_\infty^{pq}"'] & F^pH^{p+q}/F^{p+1}H^{p+q} \arrow[d,"\ol u^{p+q}"] \\
  E'^{pq}_\infty \arrow[r,"\beta'^{pq}"'] & F^pH'^{p+q}/F^{p+1}H'^{p+q}
  \end{tikzcd}\]
\end{enumerate}

我们用显然的方式定义谱序列态射的复合。有双边逆的谱序列态射称为谱序列的\emph{同构}。

若谱序列 $E$ 具有以下性质，则称之为\emph{双正则}\idx{shuangzhengze}{双正则}的：
\begin{enumerate}[label=(\roman*),leftmargin=2.8em]
  \item 对每一对 $(p,q)$，当 $k$ 充分大时有
  \[
  B_\infty(E_2^{pq})=B_k(E_2^{pq}),\qquad Z_\infty(E_2^{pq})=Z_k(E_2^{pq}).
  \]
  \item 对每个 $n$，存在 $p$ 使 $F^pH^n=0$，并存在 $p'$ 使 $F^{p'}H^n=H^n$。
\end{enumerate}

若 $\cC$ 中的谱序列 $E$ 具有以下性质，则称之为\emph{正则}\idx{zhengze}{正则}的：
\begin{enumerate}[label=(\roman*),leftmargin=2.8em]
  \item 对每一对 $(p,q)$，有
  \[
  B_\infty(E_2^{pq})=\dlim_{2\le k<\infty}B_k(E_2^{pq}),\qquad Z_\infty(E_2^{pq})=Z_k(E_2^{pq}),
  \]
  其中第二个等式对充分大的 $k$ 成立。
  \item 对每个 $n$，存在 $p$ 使 $F^pH^n=0$，且 $H^n=\dlim_pF^pH^n$。
\end{enumerate}
\begin{tnote}
(i) 第二式 $Z_\infty(E_2^{pq})=Z_k(E_2^{pq})$ 是对充分大的 $k$ 而言：在（双）正则谱序列中，固定 $(p,q)$ 后 $Z_k$ 与 $B_k$ 当 $k$ 足够大时不再变化，故 $Z_\infty$ 就等于某个 $Z_k$；第一式则把 $B_\infty$ 写成递增并 $\bigcup_kB_k$ 的有向极限。两种写法说的是同一件事。[依据：原书 p.151 (i) 与本节谱序列定义 (b)]
\end{tnote}

以下我们只对双正则谱序列陈述结论。这些结论稍作修改后对正则谱序列也成立。

\begin{proposition}\label{it:2.2.1}
设 $u\colon E\to E'$ 是双正则谱序列的态射。若对一切 $p,q\in\bZ$，$u_2^{pq}$ 都是同构，则 $u$ 是同构。
\end{proposition}
\begin{proof}
因对一切 $p,q\in\bZ$，$u_2^{pq}$ 是同构，故对一切 $p,q,r\in\bZ$（$r\ge 2$），$u_r^{pq}$ 是同构。于是对一切 $k\ge 2$ 有
\[
u_2^{pq}(B_k(E_2^{pq}))=B_k(E'^{pq}_2),\qquad u_2^{pq}(Z_k(E_2^{pq}))=Z_k(E'^{pq}_2).
\]
由于 $E$ 与 $E'$ 都是双正则的，上式推出
\[
u_2^{pq}(B_\infty(E_2^{pq}))=B_\infty(E'^{pq}_2),\qquad u_2^{pq}(Z_\infty(E_2^{pq}))=Z_\infty(E'^{pq}_2).
\]
所以对一切 $p,q\in\bZ$，$u_\infty^{pq}$ 是同构。因此对一切 $p,n\in\bZ$，
\[
\ol u^{p,n-p}\colon F^pH^n/F^{p+1}H^n\to F^pH'^n/F^{p+1}H'^n
\]
是同构。

我们来证明对一切 $n\in\bZ$，$u^n\colon H^n\to H'^n$ 是同构。由于 $E$ 与 $E'$ 都是双正则的，只需证明 $u^n$ 对一切 $p$ 诱导同构 $F^pH^n\to F^pH'^n$。我们对 $p$ 作递降归纳。当 $p$ 充分大时有 $F^pH^n=F^pH'^p=0$，此时 $F^pH^n\to F^pH'^n$ 显然是同构。
\begin{tnote}
此处原书印作「$F^pH^n=F^pH'^p=0$」，第二个等式中的上标 $p$ 应为 $n$：$F^p$ 依赖于 $H^n$ 与 $H'^n$，与 $p$ 无关的指标只有 $n$。[依据：原书 p.151；与本页上文 $u^n\colon H^n\to H'^n$ 的记号对照]
\end{tnote}
设已证 $u$ 诱导同构 $F^pH^n\to F^pH'^n$。考虑交换图
\[\begin{tikzcd}[column sep=2.4em,row sep=2.4em]
0 \arrow[r] & F^pH^n \arrow[r] \arrow[d] & F^{p-1}H^n \arrow[r] \arrow[d] & F^{p-1}H^n/F^pH^n \arrow[r] \arrow[d] & 0 \\
0 \arrow[r] & F^pH'^n \arrow[r] & F^{p-1}H'^n \arrow[r] & F^{p-1}H'^n/F^pH'^n \arrow[r] & 0
\end{tikzcd}\]
由归纳假设，第一个竖箭头是同构。我们已看到第三个竖箭头是同构。故第二个竖箭头也是同构。这就证明了断言。
\end{proof}

\begin{proposition}\label{it:2.2.2}
设 $E$ 是双正则谱序列。若存在 $r\ge 2$ 与 $q_0\in\bZ$（resp.\ $p_0\in\bZ$），使得当 $q\ne q_0$（resp.\ $p\ne p_0$）时 $E_r^{pq}=0$，则对一切 $n\in\bZ$ 有
\[
E_r^{n-q_0,q_0}\cong H^n\qquad(\text{resp.}\ E_r^{p_0,n-p_0}\cong H^n).
\]
\end{proposition}

满足命题中条件的谱序列称为\emph{退化}\idx{tuihua}{退化}的。

\begin{proof}
设当 $q\ne q_0$ 时 $E_r^{pq}=0$。则当 $q\ne q_0$ 且 $k\ge r$ 时 $E_k^{pq}=0$。于是对一切 $p,q\in\bZ$ 与 $k\ge r$，$d_k^{pq}=0$。这说明
\[
\begin{aligned}
B_r(E_2^{pq})&=B_{r+1}(E_2^{pq})=\cdots,\\
\cdots&=Z_{r+1}(E_2^{pq})=Z_r(E_2^{pq}).
\end{aligned}
\]
由于 $E$ 是双正则的，我们有
\[
\begin{aligned}
B_r(E_2^{pq})&=B_{r+1}(E_2^{pq})=\cdots=B_\infty(E_2^{pq}),\\
Z_\infty(E_2^{pq})&=\cdots=Z_{r+1}(E_2^{pq})=Z_r(E_2^{pq}).
\end{aligned}
\]
于是
\[
E_\infty^{pq}=Z_\infty(E_2^{pq})/B_\infty(E_2^{pq})=Z_r(E_2^{pq})/B_r(E_2^{pq})\cong E_r^{pq}.
\]
我们有
\[
F^pH^n/F^{p+1}H^n\cong E_\infty^{p,n-p}.
\]
所以当 $p\ne n-q_0$ 时 $F^pH^n/F^{p+1}H^n=0$，从而
\[
\cdots=F^{n-q_0-1}H^n=F^{n-q_0}H^n\supset F^{n-q_0+1}H^n=F^{n-q_0+2}H^n=\cdots.
\]
由于 $E$ 是正则的，这说明
\[
F^{n-q_0}H^n=H^n,\qquad F^{n-q_0+1}H^n=0.
\]
于是
\[
H^n\cong F^{n-q_0}H^n/F^{n-q_0+1}H^n\cong E_\infty^{n-q_0,q_0}\cong E_r^{n-q_0,q_0}.
\]
\end{proof}

\begin{proposition}\label{it:2.2.3}
设 $E$ 是双正则谱序列。若当 $p<0$、或 $q<0$、或 $0<q<n$ 时 $E_2^{pq}=0$，则对一切 $i<n$ 有 $E_2^{i,0}\cong H^i$，且有正合序列
\[
0\to E_2^{n,0}\to H^n\to E_2^{0,n}\to E_2^{n+1,0}\to H^{n+1}.
\]
特别地，若当 $p<0$ 或 $q<0$ 时 $E_2^{pq}=0$，则 $E_2^{0,0}\cong H^0$，且有正合序列
\[
0\to E_2^{1,0}\to H^1\to E_2^{0,1}\to E_2^{2,0}\to H^2.
\]
\end{proposition}
\begin{proof}
当 $p<0$、或 $q<0$、或 $0<q<n$ 时，对一切 $2\le r\le\infty$ 有 $E_r^{pq}=0$。设 $i<n$。则当 $p\ne i$ 时有
\[
F^pH^i/F^{p+1}H^i\cong E_\infty^{p,i-p}=0,
\]
从而
\[
\cdots=F^{i-1}H^i=F^iH^i\supset F^{i+1}H^i=F^{i+2}H^i=\cdots.
\]
由于 $E$ 是双正则的，这说明 $F^iH^i=H^i$ 且 $F^{i+1}H^i=0$，于是 $H^i\cong E_\infty^{i,0}$。考虑态射
\[
E_r^{i-r,r-1}\xrightarrow{d_r^{i-r,r-1}}E_r^{i,0}\xrightarrow{d_r^{i,0}}E_r^{i+r,-r+1}.
\]
因 $-r+1<0$，有 $E_r^{i+r,-r+1}=0$。因 $i-r<0$ 或 $0<r-1<n$，有 $E_r^{i-r,r-1}=0$。所以 $d_r^{i-r,r-1}$ 与 $d_r^{i,0}$ 都为零。于是
\[
\begin{aligned}
B_2(E_2^{i,0})&=B_3(E_2^{i,0})=\cdots,\\
\cdots&=Z_3(E_2^{i,0})=Z_2(E_2^{i,0}).
\end{aligned}
\]
由于 $E$ 是双正则的，这说明
\[
\begin{aligned}
B_2(E_2^{i,0})&=B_3(E_2^{i,0})=\cdots=B_\infty(E_2^{i,0}),\\
Z_\infty(E_2^{i,0})&=\cdots=Z_3(E_2^{i,0})=Z_2(E_2^{i,0}).
\end{aligned}
\]
故 $E_\infty^{i,0}\cong E_2^{i,0}$，从而对每个 $i<n$ 有 $H^i\cong E_2^{i,0}$。

当 $p\ne 0,n$ 时有 $F^pH^n/F^{p+1}H^n\cong E_\infty^{p,n-p}=0$。于是
\[
\cdots=F^{-1}H^n=F^0H^n\supset F^{-1}H^n=\cdots=F^nH^n\supset F^{n+1}H^n=\cdots.
\]
\begin{tnote}
上式中间出现的 $F^{-1}H^n$ 应为 $F^1H^n$：按 $F^\bullet$ 的递降性，$\supset$ 右边应依次为 $F^1,F^2,\dots$，且此处只有 $p=0,n$ 两个指标使商不为零，故 $F^1H^n=\cdots=F^nH^n$。[依据：原书 p.153；与同页上文 $F^pH^n/F^{p+1}H^n=0\ (p\ne 0,n)$ 对照可得]
\end{tnote}
由于 $E$ 是双正则的，这说明 $H^n=F^0H^n$ 且 $F^{n+1}H^n=0$。于是正合序列
\[
0\to F^nH^n\to H^n\to H^n/F^nH^n\to 0
\]
可等同于
\begin{equation*}\tag{1}
0\to E_\infty^{n,0}\to H^n\to E_\infty^{0,n}\to 0.
\end{equation*}
当 $p\ge n+2$ 时有 $F^pH^{n+1}/F^{p+1}H^{n+1}\cong E_\infty^{p,n+1-p}=0$。于是
\[
F^{n+2}H^{n+1}=F^{n+3}H^{n+1}=\cdots.
\]
由于 $E$ 是双正则的，必有 $F^{n+2}H^{n+1}=0$。于是正合序列
\[
0\to F^{n+1}H^{n+1}\to H^{n+1}
\]
可等同于
\begin{equation*}\tag{2}
0\to E_\infty^{n+1,0}\to H^{n+1}.
\end{equation*}

考虑态射
\[
E_r^{n-r,r-1}\xrightarrow{d_r^{n-r,r-1}}E_r^{n,0}\xrightarrow{d_r^{n,0}}E_r^{n+r,-r+1}.
\]
因 $n-r<0$ 或 $0<r-1<n$，有 $E_r^{n-r,r-1}=0$。因 $-r+1<0$，有 $E_r^{n+r,-r+1}=0$。所以 $d_r^{n-r,r-1}$ 与 $d_r^{n,0}$ 都为零。结合 $E$ 的双正则性可知
\[
\begin{aligned}
B_2(E_2^{n,0})&=B_3(E_2^{n,0})=\cdots=B_\infty(E_2^{n,0}),\\
Z_\infty(E_2^{n,0})&=\cdots=Z_3(E_2^{n,0})=Z_2(E_2^{n,0}).
\end{aligned}
\]
于是
\begin{equation*}\tag{3}
E_\infty^{n,0}\cong E_2^{n,0}.
\end{equation*}

考虑态射
\[
E_r^{-r,n+r-1}\xrightarrow{d_r^{-r,n+r-1}}E_r^{0,n}\xrightarrow{d_r^{0,n}}E_r^{r,n-r+1}.
\]
因 $-r<0$，有 $E_r^{-r,n+r-1}=0$。当 $r\ne n+1$ 时，因 $n-r+1<0$ 或 $0<n-r+1<n$，有 $E_r^{r,n-r+1}=0$。所以 $d_r^{-r,n+r-1}=0$；且当 $r\ne n+1$ 时 $d_r^{0,n}=0$。结合 $E$ 的双正则性可得
\[
\begin{aligned}
B_2(E_2^{0,n})&=\cdots=B_\infty(E_2^{0,n}),\\
Z_\infty(E_2^{0,n})&=\cdots=Z_{n+2}(E_2^{0,n})\subset Z_{n+1}(E_2^{0,n})=\cdots=Z_2(E_2^{0,n}).
\end{aligned}
\]
于是
\begin{equation*}\tag{4}
E_\infty^{0,n}\cong Z_{n+2}(E_2^{0,n})/B_{n+1}(E_2^{0,n}),
\end{equation*}
\begin{equation*}\tag{5}
E_2^{0,n}\cong E_{n+1}^{0,n}.
\end{equation*}

考虑态射
\[
E_r^{n+1-r,r-1}\xrightarrow{d_r^{n+1-r,r-1}}E_r^{n+1,0}\xrightarrow{d_r^{n+1,0}}E_r^{n+1+r,-r+1}.
\]
因 $-r+1<0$，有 $E_r^{n+1+r,-r+1}=0$。当 $r\ne n+1$ 时，因 $n+1-r<0$ 或 $0<r-1<n$，有 $E_r^{n+1-r,r-1}=0$。所以 $d_r^{n+1,0}=0$；且当 $r\ne n+1$ 时 $d_r^{n+1-r,r-1}=0$。结合 $E$ 的双正则性可得
\[
\begin{aligned}
B_2(E_2^{n+1,0})&=\cdots=B_{n+1}(E_2^{n+1,0})\subset B_{n+2}(E_2^{n+1,0})=\cdots=B_\infty(E_2^{n+1,0}),\\
Z_\infty(E_2^{n+1,0})&=\cdots=Z_2(E_2^{n+1,0}).
\end{aligned}
\]
于是
\begin{equation*}\tag{6}
E_\infty^{n+1,0}\cong Z_{n+1}(E_2^{n+1,0})/B_{n+2}(E_2^{n+1,0}),
\end{equation*}
\begin{equation*}\tag{7}
E_2^{n+1,0}\cong E_{n+1}^{n+1,0}.
\end{equation*}

通过等同
\[
\begin{aligned}
E_{n+1}^{0,n}&\cong Z_{n+1}(E_2^{0,n})/B_{n+1}(E_2^{0,n}),\\
E_{n+1}^{n+1,0}&\cong Z_{n+1}(E_2^{n+1,0})/B_{n+1}(E_2^{n+1,0}),
\end{aligned}
\]
态射
\[
d_{n+1}^{0,n}\colon E_{n+1}^{0,n}\to E_{n+1}^{n+1,0}
\]
的核等同于 $Z_{n+2}(E_2^{0,n})/B_{n+1}(E_2^{0,n})$，其余核等同于 $Z_{n+1}(E_2^{n+1,0})/B_{n+2}(E_2^{n+1,0})$。计及同构 (4)–(7)，我们得到态射
\[
E_2^{0,n}\to E_2^{n+1,0},
\]
它由 $d_{n+1}^{0,n}$ 诱导，其核同构于 $E_\infty^{0,n}$，其余核同构于 $E_\infty^{n+1,0}$。结合正合序列 (1)、(2) 与同构 (3)，我们得到正合序列
\[
0\to E_2^{n,0}\to H^n\to E_2^{0,n}\to E_2^{n+1,0}\to H^{n+1}.
\]
\end{proof}

\begin{proposition}\label{it:2.2.4}
设 $E$ 是双正则谱序列，且对一切 $p\ne p_1,p_2$（$p_1<p_2$）有 $E_2^{pq}=0$。若 $p_2-p_1=1$，则对一切 $n$ 有短正合序列
\[
0\to E_2^{p_2,n-p_2}\to H^n\to E_2^{p_1,n-p_1}\to 0.
\]
若 $p_2-p_1\ge 2$，则有长正合序列
\[
\cdots\to E_2^{p_2,n-p_2}\to H^n\to E_2^{p_1,n-p_1}\to E_2^{p_2,n+1-p_2}\to H^{n+1}\to\cdots.
\]
\end{proposition}
\begin{proof}
当 $2\le r\le\infty$ 且 $p\ne p_1,p_2$ 时 $E_r^{pq}=0$。于是当 $p\ne p_1,p_2$ 时 $F^pH^n/F^{p+1}H^n=0$，从而
\[
\cdots=F^{p_1}H^n\supset F^{p_1+1}H^n=\cdots=F^{p_2}H^n\supset F^{p_2+1}H^n=\cdots.
\]
由于 $E$ 是双正则的，这说明 $F^{p_1}H^n=H^n$ 且 $F^{p_2+1}H^n=0$。正合序列
\[
0\to F^{p_2}H^n\to H^n\to H^n/F^{p_2}H^n\to 0
\]
可等同于
\begin{equation*}\tag{1}
0\to E_\infty^{p_2,n-p_2}\to H^n\to E_\infty^{p_1,n-p_1}\to 0.
\end{equation*}
设 $p_2-p_1=1$。考虑态射
\[
E_r^{p_1-r,n-p_1+r-1}\xrightarrow{d_r^{p_1-r,n-p_1+r-1}}E_r^{p_1,n-p_1}\xrightarrow{d_r^{p_1,n-p_1}}E_r^{p_1+r,n-p_1-r+1}
\]
以及态射
\[
E_r^{p_2-r,n-p_2+r-1}\xrightarrow{d_r^{p_2-r,n-p_2+r-1}}E_r^{p_2,n-p_2}\xrightarrow{d_r^{p_2,n-p_2}}E_r^{p_2+r,n-p_2-r+1}.
\]
由于 $p_1-r$、$p_1+r$、$p_2-r$ 与 $p_2+r$ 都不同于 $p_1$ 与 $p_2$，$E_r^{p_1-r,n-p_1+r-1}$、$E_r^{p_1+r,n-p_1-r+1}$、$E_r^{p_2-r,n-p_2+r-1}$ 与 $E_r^{p_2+r,n-p_2-r+1}$ 全都为零。所以 $d_r^{p_1-r,n-p_1+r-1}$、$d_r^{p_1,n-p_1}$、$d_r^{p_2-r,n-p_2+r-1}$ 与 $d_r^{p_2,n-p_2}$ 全都为零。结合 $E$ 的双正则性可知 $E_\infty^{p_2,n-p_2}\cong E_2^{p_2,n-p_2}$ 且 $E_\infty^{p_1,n-p_1}\cong E_2^{p_1,n-p_1}$。于是正合序列 (1) 可等同于
\[
0\to E_2^{p_2,n-p_2}\to H^n\to E_2^{p_1,n-p_1}\to 0.
\]

设 $p_2-p_1\ge 2$。考虑态射
\[
E_r^{p_1-r,n-p_1+r-1}\xrightarrow{d_r^{p_1-r,n-p_1+r-1}}E_r^{p_1,n-p_1}\xrightarrow{d_r^{p_1,n-p_1}}E_r^{p_1+r,n-p_1-r+1}.
\]
有 $E_r^{p_1-r,n-p_1+r-1}=0$，且当 $r\ne p_2-p_1$ 时有 $E_r^{p_1+r,n-p_1-r+1}=0$。结合 $E$ 的双正则性可得
\[
\begin{aligned}
B_2(E_2^{p_1,n-p_1})&=\cdots=B_\infty(E_2^{p_1,n-p_1}),\\
Z_\infty(E_2^{p_1,n-p_1})&=\cdots=Z_{p_2-p_1+1}(E_2^{p_1,n-p_1})\subset\\
&\subset Z_{p_2-p_1}(E_2^{p_1,n-p_1})=\cdots=Z_2(E_2^{p_1,n-p_1}).
\end{aligned}
\]
于是
\begin{equation*}\tag{2}
E_\infty^{p_1,n-p_1}\cong Z_{p_2-p_1+1}(E_2^{p_1,n-p_1})/B_{p_2-p_1}(E_2^{p_1,n-p_1}),
\end{equation*}
\begin{equation*}\tag{3}
E_2^{p_1,n-p_1}\cong E_{p_2-p_1}^{p_1,n-p_1}.
\end{equation*}

考虑态射
\[
E_r^{p_2-r,n-p_2+r}\xrightarrow{d_r^{p_2-r,n-p_2+r}}E_r^{p_2,n+1-p_2}\xrightarrow{d_r^{p_2,n+1-p_2}}E_r^{p_2+r,n+2-p_2-r}.
\]
有 $E_r^{p_2+r,n+2-p_2-r}=0$，且当 $r\ne p_2-p_1$ 时有 $E_r^{p_2-r,n-p_2+r}=0$。结合 $E$ 的双正则性可得
\[
\begin{aligned}
B_2(E_2^{p_2,n+1-p_2})&=\cdots=B_{p_2-p_1}(E_2^{p_2,n+1-p_2})\subset\\
&\subset B_{p_2-p_1+1}(E_2^{p_2,n+1-p_2})=\cdots=B_\infty(E_2^{p_2,n+1-p_2}),\\
Z_\infty(E_2^{p_2,n+1-p_2})&=\cdots=Z_2(E_2^{p_2,n+1-p_2}).
\end{aligned}
\]
于是
\begin{equation*}\tag{4}
E_\infty^{p_2,n+1-p_2}\cong Z_{p_2-p_1}(E_2^{p_2,n+1-p_2})/B_{p_2-p_1+1}(E_2^{p_2,n+1-p_2}),
\end{equation*}
\begin{equation*}\tag{5}
E_2^{p_2,n+1-p_2}\cong E_{p_2-p_1}^{p_2,n+1-p_2}.
\end{equation*}

通过等同
\[
\begin{aligned}
E_{p_2-p_1}^{p_1,n-p_1}&\cong Z_{p_2-p_1}(E_2^{p_1,n-p_1})/B_{p_2-p_1}(E_2^{p_1,n-p_1}),\\
E_{p_2-p_1}^{p_2,n+1-p_2}&\cong Z_{p_2-p_1}(E_2^{p_2,n+1-p_2})/B_{p_2-p_1}(E_2^{p_2,n+1-p_2}),
\end{aligned}
\]
态射
\[
d_{p_2-p_1}^{p_1,n-p_1}\colon E_{p_2-p_1}^{p_1,n-p_1}\to E_{p_2-p_1}^{p_2,n+1-p_2}
\]
的核等同于
\[
Z_{p_2-p_1+1}(E_2^{p_1,n-p_1})/B_{p_2-p_1}(E_2^{p_1,n-p_1}),
\]
其余核等同于
\[
Z_{p_2-p_1}(E_2^{p_2,n+1-p_2})/B_{p_2-p_1+1}(E_2^{p_2,n+1-p_2}).
\]
计及同构 (2)–(5)，我们得到态射
\[
E_2^{p_1,n-p_1}\to E_2^{p_2,n+1-p_2},
\]
它由 $d_{p_2-p_1}^{p_1,n-p_1}$ 诱导，其核等同于 $E_\infty^{p_1,n-p_1}$，其余核等同于 $E_\infty^{p_2,n+1-p_2}$。结合正合序列 (1)，我们得到长正合序列
\[
\cdots\to E_2^{p_2,n-p_2}\to H^n\to E_2^{p_1,n-p_1}\to E_2^{p_2,n+1-p_2}\to H^{n+1}\to\cdots.
\]
\end{proof}

下面命题的证明留给读者：

\begin{proposition}\label{it:2.2.5}
设 $E$ 是双正则谱序列，且对一切 $q\ne q_1,q_2$（$q_1<q_2$）有 $E_2^{pq}=0$。则有长正合序列
\[
\cdots\to E_2^{n-q_1,q_1}\to H^n\to E_2^{n-q_2,q_2}\to E_2^{n+1-q_1,q_1}\to H^{n+1}\to\cdots.
\]
\end{proposition}

下面的引理是显然的：

\begin{lemma}\label{it:2.2.6}
设给定交换图
\[\begin{tikzcd}[column sep=2.6em,row sep=2.0em]
& D \arrow[d,"\gamma"] & \\
A \arrow[ur] \arrow[r,"\alpha"'] & B \arrow[r,"\beta"'] & C
\end{tikzcd}\]
且下排正合。则 $\beta$ 诱导同构
\[
\Image\gamma/\Image\alpha\to\Image(\beta\gamma).
\]
\end{lemma}

设 $K^{\cdot}$ 是阿贝尔范畴 $\cC$ 中带有\emph{递降滤过}的\emph{过滤复形}\idx{guolvfuxing}{过滤复形}：
\[
\cdots\supset F^pK^{\cdot}\supset F^{p+1}K^{\cdot}\supset\cdots.
\]
我们将由这个过滤复形构造一个谱序列。

对任意 $p,q,r\in\bZ$ 与 $r\ge 2$，令
\[
Z_r^{pq}=\Image\bigl(H^{p+q}(F^pK^{\cdot}/F^{p+r}K^{\cdot})\to H^{p+q}(F^pK^{\cdot}/F^{p+1}K^{\cdot})\bigr),
\]
其中态射 $H^{p+q}(F^pK^{\cdot}/F^{p+r}K^{\cdot})\to H^{p+q}(F^pK^{\cdot}/F^{p+1}K^{\cdot})$ 是短正合序列
\[
0\to F^{p+1}K^{\cdot}/F^{p+r}K^{\cdot}\to F^pK^{\cdot}/F^{p+r}K^{\cdot}\to F^pK^{\cdot}/F^{p+1}K^{\cdot}\to 0
\]
所关联的上同调对象的长正合序列中的态射。令
\[
B_r^{pq}=\Image\bigl(H^{p+q-1}(F^{p-r+1}K^{\cdot}/F^pK^{\cdot})\to H^{p+q}(F^pK^{\cdot}/F^{p+1}K^{\cdot})\bigr),
\]
其中态射 $H^{p+q-1}(F^{p-r+1}K^{\cdot}/F^pK^{\cdot})\to H^{p+q}(F^pK^{\cdot}/F^{p+1}K^{\cdot})$ 是短正合序列
\[
0\to F^pK^{\cdot}/F^{p+1}K^{\cdot}\to F^{p-r+1}K^{\cdot}/F^{p+1}K^{\cdot}\to F^{p-r+1}K^{\cdot}/F^pK^{\cdot}\to 0
\]
所关联的上同调对象的长正合序列中的态射。我们约定 $F^\infty K^{\cdot}=0$、$F^{-\infty}K^{\cdot}=K^{\cdot}$，并用同样的公式定义 $Z_\infty^{pq}$ 与 $B_\infty^{pq}$。于是
\[
Z_\infty^{pq}=\Image\bigl(H^{p+q}(F^pK^{\cdot})\to H^{p+q}(F^pK^{\cdot}/F^{p+1}K^{\cdot})\bigr),
\]
\[
B_\infty^{pq}=\Image\bigl(H^{p+q-1}(K^{\cdot}/F^pK^{\cdot})\to H^{p+q}(F^pK^{\cdot}/F^{p+1}K^{\cdot})\bigr).
\]
我们有
\[
B_2^{pq}\subset B_3^{pq}\subset\cdots\subset B_\infty^{pq},
\]
\[
Z_\infty^{pq}\subset\cdots\subset Z_3^{pq}\subset Z_2^{pq}.
\]
由于有交换图
\[\begin{tikzcd}[column sep=2.2em,row sep=1.9em]
0 \arrow[r] & F^pK^{\cdot} \arrow[r] \arrow[d] & K^{\cdot} \arrow[r] \arrow[d] & K^{\cdot}/F^pK^{\cdot} \arrow[r] \arrow[d,equal] & 0\\
0 \arrow[r] & F^pK^{\cdot}/F^{p+1}K^{\cdot} \arrow[r] & K^{\cdot}/F^{p+1}K^{\cdot} \arrow[r] & K^{\cdot}/F^pK^{\cdot} \arrow[r] & 0,
\end{tikzcd}\]
所以下述图表交换：
\[\begin{tikzcd}[column sep=3.2em,row sep=1.9em]
H^{p+q-1}(K^{\cdot}/F^pK^{\cdot}) \arrow[r] \arrow[d,equal] & H^{p+q}(F^pK^{\cdot}) \arrow[d] \\
H^{p+q-1}(K^{\cdot}/F^pK^{\cdot}) \arrow[r] & H^{p+q}(F^pK^{\cdot}/F^{p+1}K^{\cdot}).
\end{tikzcd}\]
于是
\[
B_\infty^{pq}\subset Z_\infty^{pq}.
\]
定义
\[
E_r^{pq}=Z_r^{pq}/B_r^{pq}\quad(r\ge 2),\qquad E_\infty^{pq}=Z_\infty^{pq}/B_\infty^{pq}.
\]
把 \ref{it:2.2.6} 应用于交换图
\[\begin{tikzcd}[column sep=2.4em,row sep=1.7em,cells={nodes={font=\small}}]
& H^{p+q}(F^pK^{\cdot}/F^{p+r}K^{\cdot}) \arrow[dl] \arrow[d] & \\
H^{p+q}(F^pK^{\cdot}/F^{p+r+1}K^{\cdot}) \arrow[r] & H^{p+q}(F^pK^{\cdot}/F^{p+1}K^{\cdot}) \arrow[r] & H^{p+q+1}(F^{p+1}K^{\cdot}/F^{p+r+1}K^{\cdot}),
\end{tikzcd}\]
我们看到 $Z_r^{pq}/Z_{r+1}^{pq}$ 同构于复合
\[
H^{p+q}(F^pK^{\cdot}/F^{p+r}K^{\cdot})\to H^{p+q}(F^pK^{\cdot}/F^{p+1}K^{\cdot})\to H^{p+q+1}(F^{p+1}K^{\cdot}/F^{p+r+1}K^{\cdot})
\]
的像。把 \ref{it:2.2.6} 应用于交换图
\[\begin{tikzcd}[column sep=2.4em,row sep=1.7em,cells={nodes={font=\small}}]
& H^{p+q}(F^pK^{\cdot}/F^{p+r}K^{\cdot}) \arrow[dl] \arrow[d] & \\
H^{p+q}(F^{p+1}K^{\cdot}/F^{p+r}K^{\cdot}) \arrow[r] & H^{p+q+1}(F^{p+r}K^{\cdot}/F^{p+r+1}K^{\cdot}) \arrow[r] & H^{p+q+1}(F^{p+1}K^{\cdot}/F^{p+r+1}K^{\cdot}),
\end{tikzcd}\]
我们看到 $B_{r+1}^{p+r,q-r+1}/B_r^{p+r,q-r+1}$ 同构于复合
\[
H^{p+q}(F^pK^{\cdot}/F^{p+r}K^{\cdot})\to H^{p+q+1}(F^{p+r}K^{\cdot}/F^{p+r+1}K^{\cdot})\to H^{p+q+1}(F^{p+1}K^{\cdot}/F^{p+r+1}K^{\cdot})
\]
的像。由于图表
\[\begin{tikzcd}[column sep=2.4em,row sep=1.9em,cells={nodes={font=\small}}]
H^{p+q}(F^pK^{\cdot}/F^{p+r}K^{\cdot}) \arrow[r] \arrow[d] & H^{p+q+1}(F^{p+r}K^{\cdot}/F^{p+r+1}K^{\cdot}) \arrow[d] \\
H^{p+q}(F^pK^{\cdot}/F^{p+1}K^{\cdot}) \arrow[r] & H^{p+q+1}(F^{p+1}K^{\cdot}/F^{p+r+1}K^{\cdot})
\end{tikzcd}\]
交换，我们有
\[
Z_r^{pq}/Z_{r+1}^{pq}\cong B_{r+1}^{p+r,q-r+1}/B_r^{p+r,q-r+1}.
\]
定义
\[
d_r^{pq}\colon E_r^{pq}\to E_r^{p+r,q-r+1}
\]
为复合
\[\begin{tikzcd}[column sep=2.0em,row sep=1.6em]
Z_r^{pq}/B_r^{pq} \arrow[r] & Z_r^{pq}/Z_{r+1}^{pq} \arrow[r,"\cong"] & B_{r+1}^{p+r,q-r+1}/B_r^{p+r,q-r+1} \arrow[r] & Z_r^{p+r,q-r+1}/B_r^{p+r,q-r+1}.
\end{tikzcd}\]
我们有
\[
\Ker d_r^{pq}=Z_{r+1}^{pq}/B_r^{pq},\qquad \Image d_r^{pq}=B_{r+1}^{p+r,q-r+1}/B_r^{p+r,q-r+1}.
\]
这说明
\[
d_r^{p+r,q-r+1}d_r^{pq}=0
\]
且
\[
\Ker d_r^{pq}/\Image d_r^{p-r,q+r-1}\cong Z_{r+1}^{pq}/B_{r+1}^{pq}=E_{r+1}^{pq}.
\]
对每个 $n\in\bZ$，定义 $H^n=H^n(K^{\cdot})$，并用
\[
F^pH^n=\Image\bigl(H^n(F^pK^{\cdot})\to H^n(K^{\cdot})\bigr)
\]
在 $H^n$ 上定义递降滤过。把 \ref{it:2.2.6} 应用于图表
\[\begin{tikzcd}[column sep=2.4em,row sep=1.7em,cells={nodes={font=\small}}]
& H^{p+q}(F^pK^{\cdot}) \arrow[dl] \arrow[d] & \\
H^{p+q}(F^{p+1}K^{\cdot}) \arrow[r] & H^{p+q}(K^{\cdot}) \arrow[r] & H^{p+q}(K^{\cdot}/F^{p+1}K^{\cdot}),
\end{tikzcd}\]
我们看到 $F^pH^{p+q}/F^{p+1}H^{p+q}$ 同构于复合
\[
H^{p+q}(F^pK^{\cdot})\to H^{p+q}(K^{\cdot})\to H^{p+q}(K^{\cdot}/F^{p+1}K^{\cdot})
\]
的像。把 \ref{it:2.2.6} 应用于图表
\[\begin{tikzcd}[column sep=2.4em,row sep=1.7em,cells={nodes={font=\small}}]
& H^{p+q}(F^pK^{\cdot}) \arrow[dl] \arrow[d] & \\
H^{p+q-1}(K^{\cdot}/F^pK^{\cdot}) \arrow[r] & H^{p+q}(F^pK^{\cdot}/F^{p+1}K^{\cdot}) \arrow[r] & H^{p+q}(K^{\cdot}/F^{p+1}K^{\cdot}),
\end{tikzcd}\]
我们看到 $Z_\infty^{pq}/B_\infty^{pq}$ 同构于复合
\[
H^{p+q}(F^pK^{\cdot})\to H^{p+q}(F^pK^{\cdot}/F^{p+1}K^{\cdot})\to H^{p+q}(K^{\cdot}/F^{p+1}K^{\cdot})
\]
的像。由于图表
\[\begin{tikzcd}[column sep=2.4em,row sep=1.9em,cells={nodes={font=\small}}]
H^{p+q}(F^pK^{\cdot}) \arrow[r] \arrow[d] & H^{p+q}(K^{\cdot}) \arrow[d] \\
H^{p+q}(F^pK^{\cdot}/F^{p+1}K^{\cdot}) \arrow[r] & H^{p+q}(K^{\cdot}/F^{p+1}K^{\cdot})
\end{tikzcd}\]
交换，我们有
\[
F^pH^{p+q}/F^{p+1}H^{p+q}\cong Z_\infty^{pq}/B_\infty^{pq}.
\]
最后我们定义 $B_\infty(E_2^{pq})$ 与 $Z_\infty(E_2^{pq})$ 为 $E_2^{pq}$ 的子对象 $B_\infty^{pq}/B_2^{pq}$ 与 $Z_\infty^{pq}/B_2^{pq}$。至此谱序列的构造完成。

注意，若定义
\[
E_1^{pq}=H^{p+q}(F^pK^{\cdot}/F^{p+1}K^{\cdot}),
\]
并定义 $d_1^{pq}\colon E_1^{pq}\to E_1^{p+1,q}$ 为短正合序列
\[
0\to F^{p+1}K^{\cdot}/F^{p+2}K^{\cdot}\to F^pK^{\cdot}/F^{p+2}K^{\cdot}\to F^pK^{\cdot}/F^{p+1}K^{\cdot}\to 0
\]
所关联的上同调群的长正合序列中的态射
\[
H^{p+q}(F^pK^{\cdot}/F^{p+1}K^{\cdot})\to H^{p+q+1}(F^{p+1}K^{\cdot}/F^{p+2}K^{\cdot}),
\]
则
\[
E_2^{pq}=\Ker d_1^{pq}/\Image d_1^{p-1,q}.
\]
我们常把 $(E_1^{pq},d_1^{pq})$ 看作上述谱序列的一部分，并把该谱序列记为
\[
E_1^{pq}=H^{p+q}(F^pK^{\cdot}/F^{p+1}K^{\cdot})\Rightarrow H^{p+q}(K^{\cdot}).
\]
我们还可以令
\[
E_0^{pq}=F^pK^{p+q}/F^{p+1}K^{p+q},
\]
并且对任意 $r\ge 0$，令
\[
Z_r^{pq}=\bigl(F^pK^{p+q}\cap d^{-1}(F^{p+r}K^{p+q+1})+F^{p+1}K^{p+q}\bigr)/F^{p+1}K^{p+q},
\]
\[
B_r^{pq}=\bigl(F^pK^{p+q}\cap d(F^{p-r+1}K^{p+q-1})+F^{p+1}K^{p+q}\bigr)/F^{p+1}K^{p+q},
\]
\[
E_r^{pq}=Z_r^{pq}/B_r^{pq}\cong\frac{F^pK^{p+q}\cap d^{-1}(F^{p+r}K^{p+q+1})}{F^pK^{p+q}\cap d(F^{p-r+1}K^{p+q-1})+F^{p+1}K^{p+q}\cap d^{-1}(F^{p+r}K^{p+q+1})},
\]
并定义 $d_r^{pq}\colon E_r^{pq}\to E_r^{p+r,q-r+1}$ 为复合
\begin{align*}
E_r^{pq}&\cong\frac{F^pK^{p+q}\cap d^{-1}(F^{p+r}K^{p+q+1})}{F^pK^{p+q}\cap d(F^{p-r+1}K^{p+q-1})+F^{p+1}K^{p+q}\cap d^{-1}(F^{p+r}K^{p+q+1})}\\
&\xrightarrow{\ d\ }\frac{F^{p+r}K^{p+q+1}\cap d^{-1}(F^{p+2r}K^{p+q+2})}{F^{p+r}K^{p+q+1}\cap d(F^{p+1}K^{p+q})+F^{p+r+1}K^{p+q+1}\cap d^{-1}(F^{p+2r}K^{p+q+2})}\\
&\cong E_r^{p+r,q-r+1},
\end{align*}
其中第二个态射由 $d\colon K^{p+q}\to K^{p+q+1}$ 诱导。可以验证
\[
\Ker d_r^{pq}/\Image d_r^{p-r,q+r-1}\cong E_{r+1}^{pq}.
\]
这给出了谱序列的另一种构造。细节留给读者。

